\scanpage{260}
- \mspage{99} - \margin{IX}

que l'on a une \uline{immersion ouverte} canonique
\[
(8.3.1)\qquad A_{K,K'}^{\unsure{\natural|\sharp}}
 \hookrightarrow A_{K,K''}^{\unsure{\natural|\sharp}},
\]
permettant d'identifier le premier groupe à un sous-groupe ouvert du deuxième.
Pour \(K'\) variable, on voit que les \(A_{K,K'}\) forment ainsi un système inductif
de groupes lisses de type fini sur \(K\), à morphismes de transition des immersions
ouvertes. Si on pose
\[
(8.3.2)\qquad\Phi_{K,K'}=
 A_{K,K'}^{\unsure{\natural|\sharp}}/A_{K,K'}^{\unsure{\natural|\sharp}\,0},
\]
les \(\Phi_{K,K'}\) eux-mêmes forment donc un système inductif de groupes finis
étales sur \(K\), qui sera étudié plus en détail plus bas.

Nous poserons
\[
(8.3.3)\qquad A_K^{\unsure{b|\flat}}=
 \varinjlim_{K'}A_{K,K'}^{\unsure{\natural|\sharp}},
\]
la limite étant prise dans la catégorie des faisceaux fppf sur \(K\). Comme
les morphismes de transition dans le système inductif envisagé dans des immersions
ouvertes, il est d'ailleurs clair que cette limite est en fait représentable
par un schéma en groupes (en général non quasi-compact) sur \(K\), que nous
désignerons comme de juste par le même symbole \(A_K^{\unsure{b|\flat}}\).
On notera qu'on a, pour ce schéma en groupes
\margin{\figure{\begin{tikzpicture}[scale=.3]
\draw[line width=1pt] (0,0)--(.8,0);\draw (.1,.3)--(.3,.35)--(.4,.3);
\end{tikzpicture}}}
\[
(8.3.4)\qquad A_K^{\unsure{b|\flat}\,0}=A_K^{\unsure{\natural|\sharp}\,0},\quad
 A_K^{\unsure{b|\flat}}/A_K^{\unsure{b|\flat}\,0}
 \simeq\varinjlim_{K'}\Phi_{K,K'}.
\]
Ce dernier groupe sera déterminé dans 11.10 : il est canoniquement isomorphe
à \(\underline M_K\otimes\mathbb Q/\mathbb Z\) (avec la notation (8.1.8)).
\margin{417}
\margin{417}
\margin{259}
